% Módulo central mostrando la vista de examen (finalizado).
% Se compila por separado: lualatex onx_examen_finalizado.tex -> onx_examen_finalizado.pdf
% Se dibuja sobre red_onx_apagada.png. La disposición replica screen_main.cpp.
% El panel muestra los colores invertidos: el rojo 0xe74c3c se ve celeste,
% el gris 0x888888 se ve gris medio y el texto negro se ve claro.
\documentclass[tikz,border=0pt]{standalone}
\usepackage{fontspec}
\newfontfamily\montserrat{Montserrat-Medium.otf}

\definecolor{fondo}{RGB}{42,43,46}
\definecolor{claro}{RGB}{232,240,252}
\definecolor{celeste}{RGB}{24,179,195}
\definecolor{gris}{RGB}{119,119,119}

\newcommand\fuenteveinte{\montserrat\fontsize{9.1}{10.5}\selectfont}
\newcommand\fuentedieciocho{\montserrat\fontsize{8.2}{9.6}\selectfont}
\newcommand\fuentedieciseis{\montserrat\fontsize{7.3}{8.6}\selectfont}

\begin{document}
\begin{tikzpicture}[x=0.01cm,y=0.01cm]
\node[anchor=south west, inner sep=0] at (0,0) {\includegraphics[width=6.39cm]{red_onx_apagada.png}};
% Área visible del panel en la imagen: (64,101)-(575,867), equivale a 320x480 px
\begin{scope}[shift={(64,867)}, x=0.015969cm, y=-0.015958cm]
  \fill[fondo] (0,0) rectangle (320,480);
  \begin{scope}[every node/.style={inner sep=0, text=claro, align=center}]
    \node[anchor=north, font=\fuentedieciocho] at (160,20) {EXAMEN\\FINALIZADO};
    \node[font=\fuentedieciseis, text=celeste] at (160,255) {Necesario: 3\\Obtenido: 1 / 5\\Tiempo: 01:44\\[8.6pt]DESAPROBADO};
    \node[anchor=south, font=\fuentedieciseis, text=gris] at (160,465) {Volviendo en 2 s...};
  \end{scope}
\end{scope}
\end{tikzpicture}
\end{document}
